% Part: model-theory
% Chapter: basics
% Section: nonstandard-arithmetic

\documentclass[../../../include/open-logic-section]{subfiles}

\begin{document}

\olfileid{mod}{bas}{nsa}
\olsection{في نماذج الحساب غير القياسية}

\begin{defn}
لتكن~$\Lang{L_N}$ !!{language} الحساب. وقوام هذه اللغة~!!{constant}
$\Obj{0}$، و!!{predicate} ثنائي الموضع~$<$، و!!{function} أحادي الموضع
$\prime$، و!!{function}s ثنائية الموضع~$+$ و$\times$.
\begin{enumerate}
\item أما \emph{النموذج القياسي} للحساب فهو~!!{structure}~$\Struct{N}$
  للغة~$\Lang{L_N}$، مجاله~$\Nat = \{ \OLMathNumeral{0}, \OLMathNumeral{1}, \OLMathNumeral{2}, \dots\}$. وفيه يفسر
  $\Obj{0}$ بالصفر~$\OLMathNumeral{0}$، و$<$ بعلاقة الأصغر من على~$\Nat$، وتفسر
  $\prime$ و$+$ و$\times$، على الترتيب، بالخلف والجمع والضرب على~$\Nat$.
\item وأما \emph{الحساب الصادق} فهو النظرية~$\Theory{\OLId{latin-looped}{N}}$.
\end{enumerate}
\end{defn}

إذا عملنا في~$\Lang{L_N}$ اختصرنا كل حد من الصورة
$\Obj{0}^{\prime\dots\prime}$، تكون فيه لدالة الخلف~$\OLId{latin-ordinary}{n}$ تطبيقات على
$\Obj{0}$، فرمزنا إليه بـ~$\num{\OLId{latin-ordinary}{n}}$.

\begin{defn}
ونسمي~!!{structure}~$\Struct{M}$ للغة~$\Lang{L_N}$ \emph{قياسيًا} إذا
وفقط إذا تحقق~$\Struct{N} \iso \Struct{M}$.
\end{defn}

\begin{thm}\ollabel{thm:non-std}
توجد نماذج~!!{enumerable} غير قياسية للحساب الصادق.
\end{thm}

\begin{proof}
نوسع~$\Lang{L_N}$ بـ!!{constant} جديد~$\OLId{latin-ordinary}{c}$، ثم ننظر في النظرية
\[
\Theory{\OLId{latin-looped}{N}} \cup \Setabs{\num{\OLId{latin-ordinary}{n}} < \OLId{latin-ordinary}{c}}{\OLId{latin-ordinary}{n} \in \Nat}.
\]
وهذه النظرية قابلة للإشباع منتهيًا؛ فمن التراص يكون لها نموذج~$\Struct{M}$،
وبمبرهنة لوفنهايم--سكولم النزولية يجوز أن نختاره~!!{enumerable}. فإذا كان
$\Domain{\OLId{latin-looped}{M}}$ مجال~$\Struct{M}$، أثبتنا لما يفسر به~$\Struct{M}$ الرموز
غير المنطقية في~$\Lang{L_N}$ الأسماء الآتية:
$\mathbf{z} = \Assign{\Obj{0}}{\OLId{latin-looped}{M}} \in \Domain{\OLId{latin-looped}{M}}$، و
${\prec} = \Assign{<}{\OLId{latin-looped}{M}} \subseteq \Domain{\OLId{latin-looped}{M}}^\OLMathNumeral{2}$، و
$* = \Assign{\prime}{\OLId{latin-looped}{M}} \colon \Domain{\OLId{latin-looped}{M}} \to \Domain{\OLId{latin-looped}{M}}$، و
$\oplus = \Assign{+}{\OLId{latin-looped}{M}}$، و
$\otimes = \Assign{\times}{\OLId{latin-looped}{M}} \colon \Domain{\OLId{latin-looped}{M}}^\OLMathNumeral{2} \to \Domain{\OLId{latin-looped}{M}}$.
ولكل~$\OLId{latin-ordinary}{x} \in \Domain{\OLId{latin-looped}{M}}$، نرمز بـ~$\OLId{latin-ordinary}{x}^*$ إلى العنصر من~$\Domain{\OLId{latin-looped}{M}}$
الحاصل من تطبيق~$*$ على~$\OLId{latin-ordinary}{x}$.

فلو كان~$\OLId{latin-ordinary}{h}$ تماثلًا بين~$\Struct{N}$ و$\Struct{M}$، لوجد
$\OLId{latin-ordinary}{n} \in \Nat$ بحيث~$\OLId{latin-ordinary}{h}(\OLId{latin-ordinary}{n}) = \Assign{\OLId{latin-ordinary}{c}}{\OLId{latin-looped}{M}}$. فليكن~$\OLId{latin-ordinary}{s}$ أي إسناد في
$\Struct{N}$ بحيث~$\OLId{latin-ordinary}{s}(\OLId{latin-ordinary}{x}) = \OLId{latin-ordinary}{n}$. عندئذ
$\Sat{\OLId{latin-looped}{N}}{\eq[\num{\OLId{latin-ordinary}{n}}][\OLId{latin-ordinary}{x}]}[\OLId{latin-ordinary}{s}]$؛ وببرهان~\olref[iso]{thm:isom} يكون
كذلك~$\Sat{\OLId{latin-looped}{M}}{\eq[\num{\OLId{latin-ordinary}{n}}][\OLId{latin-ordinary}{x}]}[\OLId{latin-ordinary}{h} \circ \OLId{latin-ordinary}{s}]$، ومن ثم
$\Assign{\OLId{latin-ordinary}{c}}{\OLId{latin-looped}{M}} = \mathbf{z}^{*\cdots *}$، وفي الطرف الأيمن يتكرر~$*$ عدد~$\OLId{latin-ordinary}{n}$
من المرات. وهذا محال؛ إذ يقتضي الفرض
$\Sat{\OLId{latin-looped}{M}}{\num{\OLId{latin-ordinary}{n}} < \OLId{latin-ordinary}{c}}$، و~$\prec$ لا انعكاسية. فثبت أن
$\Struct{M}$ غير قياسي.
\end{proof}

\begin{prob}
تكون علاقة~$\OLId{latin-looped}{R}$ على مجموعة~$\OLId{latin-looped}{X}$ \emph{راسخة} إذا وفقط إذا لم توجد فيها
سلاسل نزولية غير منتهية؛ أي إذا لم توجد~$\OLId{latin-ordinary}{x}_\OLMathNumeral{0}$،~$\OLId{latin-ordinary}{x}_\OLMathNumeral{1}$،~$\OLId{latin-ordinary}{x}_\OLMathNumeral{2}$،~\dots
في~$\OLId{latin-looped}{X}$ بحيث~$\dots \OLId{latin-ordinary}{x}_\OLMathNumeral{2}Rx_\OLMathNumeral{1}Rx_\OLMathNumeral{0}$. وبافتراض اتساق نظرية مجموعات
زيرميلو--فرينكل~$ZF$، بيّن أن لـ$ZF$ نماذج غير راسخة؛ أي نماذج
$\mathfrak{M}$ بحيث~$\dots \OLId{latin-ordinary}{x}_\OLMathNumeral{2} \in \OLId{latin-ordinary}{x}_\OLMathNumeral{1} \in \OLId{latin-ordinary}{x}_\OLMathNumeral{0}$.
\end{prob}

النموذج غير القياسي~$\Struct{M}$ الوارد في~\olref{thm:non-std} مكافئ
أوليًا للنموذج القياسي، ومن هذا التكافؤ تستخرج طائفة من خواص~$\Struct{M}$.
وفيما بقي من هذا القسم نفصل تلك الخواص، حتى ننتهي إلى وصف دقيق لنماذج~$\Theory{\OLId{latin-looped}{N}}$
غير القياسية~!!{enumerable}.

\begin{enumerate}
\item لا يكون أي عنصر من~$\Domain{\OLId{latin-looped}{M}}$ أصغر بحسب~$\prec$ من نفسه:
  فالجملة~$\lforall[\OLId{latin-ordinary}{x}][\lnot \OLId{latin-ordinary}{x} < \OLId{latin-ordinary}{x}]$ صادقة في~$\Struct{N}$، ومن ثم
  في~$\Struct{M}$.
\item وباستدلال مماثل ينتج لنا أن~$\prec$ \emph{ترتيب خطي} لـ
  $\Domain{\OLId{latin-looped}{M}}$، أي علاقة كلية لا انعكاسية ومتعدية على~$\Domain{\OLId{latin-looped}{M}}$.
\item العنصر~$\mathbf{z}$ هو أصغر عناصر~$\Domain{\OLId{latin-looped}{M}}$ بحسب~$\prec$.
\item كل عنصر من~$\Domain{\OLId{latin-looped}{M}}$ أصغر بحسب~$\prec$ من خلفه بالنسبة إلى
  $*$، و$\OLId{latin-ordinary}{x}^*$ هو أصغر عناصر~$\Domain{\OLId{latin-looped}{M}}$ بحسب~$\prec$ بين العناصر
  الأكبر من~$\OLId{latin-ordinary}{x}$.
\item يحتوي~$\Struct{M}$ قطعة ابتدائية بالنسبة إلى~$\prec$ متماثلة
  مع~$\Nat$، هي~$\mathbf{z}, \mathbf{z}^*, \mathbf{z}^{**}, \dots$،
  ونسميها~\emph{الجزء القياسي} من~$\Domain{\OLId{latin-looped}{M}}$. وكل عنصر آخر من
  $\Domain{\OLId{latin-looped}{M}}$ \emph{غير قياسي}. ولا بد من وجود عناصر غير قياسية في
  $\Domain{\OLId{latin-looped}{M}}$، وإلا لكانت الدالة~$\OLId{latin-ordinary}{h}$ من برهان~\olref{thm:non-std}
  تماثلًا. ونستخدم~$\OLId{latin-ordinary}{n}, \OLId{latin-ordinary}{m}, \dots$ بوصفها~!!{variable}s تتراوح على هذا
  الجزء القياسي من~$\Struct{M}$.
\item كل عنصر غير قياسي أكبر من كل عنصر قياسي؛ ذلك أنه لكل
  $\OLId{latin-ordinary}{n} \in \Nat$،
  \[
  \Sat{\OLId{latin-looped}{N}}{\lforall[\OLId{latin-ordinary}{z}][(\lnot(\eq[\OLId{latin-ordinary}{z}][\Obj{0}] \lor
  \dots \lor \eq[\OLId{latin-ordinary}{z}][\num{\OLId{latin-ordinary}{n}}]) \lif \num{\OLId{latin-ordinary}{n}} < \OLId{latin-ordinary}{z})]},
  \]
  ولذلك إذا كان~$\OLId{latin-ordinary}{z} \in \Domain{\OLId{latin-looped}{M}}$ مختلفًا عن جميع العناصر القياسية،
  وجب أن يكون~\emph{أكبر} منها جميعًا.
\item كل عنصر من~$\Domain{\OLId{latin-looped}{M}}$ غير~$\mathbf{z}$ هو خلف بالنسبة إلى~$*$
  لعنصر وحيد من~$\Domain{\OLId{latin-looped}{M}}$، نرمز إليه بـ$^*\OLId{latin-ordinary}{x}$. وإذا كان~$\OLId{latin-ordinary}{x} = \OLId{latin-ordinary}{y}^*$،
  كان~$\OLId{latin-ordinary}{x}$ و$\OLId{latin-ordinary}{y}$ قياسيين إذا كان أحدهما قياسيًا، وكانا غير قياسيين إذا
  كان أحدهما غير قياسي.
\item عرّف علاقة تكافؤ~$\approx$ على~$\Domain{\OLId{latin-looped}{M}}$ بالقول إن
  $\OLId{latin-ordinary}{x} \approx \OLId{latin-ordinary}{y}$ إذا وفقط إذا كان، لبعض~$\OLId{latin-ordinary}{n}$ \emph{قياسي}، إما
  $\OLId{latin-ordinary}{x} \oplus \OLId{latin-ordinary}{n} = \OLId{latin-ordinary}{y}$ وإما~$\OLId{latin-ordinary}{y} \oplus \OLId{latin-ordinary}{n} =\OLId{latin-ordinary}{x}$. وبعبارة أخرى، يكون
  $\OLId{latin-ordinary}{x} \approx \OLId{latin-ordinary}{y}$ إذا وفقط إذا كانت المسافة بين~$\OLId{latin-ordinary}{x}$ و$\OLId{latin-ordinary}{y}$ منتهية. وإذا
  كان~$\OLId{latin-ordinary}{n}$ و$\OLId{latin-ordinary}{m}$ قياسيين، كان~$\OLId{latin-ordinary}{n} \approx \OLId{latin-ordinary}{m}$. وعرّف~\emph{كتلة}~$\OLId{latin-ordinary}{x}$
  بأنها صنف التكافؤ~$[\OLId{latin-ordinary}{x}] = \Setabs{\OLId{latin-ordinary}{y}}{\OLId{latin-ordinary}{x} \approx \OLId{latin-ordinary}{y}}$.
\item افترض~$\OLId{latin-ordinary}{x} \prec \OLId{latin-ordinary}{y}$ و$\OLId{latin-ordinary}{x} \not\approx \OLId{latin-ordinary}{y}$. ولأن
  $\Sat{\OLId{latin-looped}{N}}{\lforall[\OLId{latin-ordinary}{x}][\lforall[\OLId{latin-ordinary}{y}][(\OLId{latin-ordinary}{x} < \OLId{latin-ordinary}{y} \lif (\OLId{latin-ordinary}{x}' < \OLId{latin-ordinary}{y} \lor \OLId{latin-ordinary}{x}' =
      \OLId{latin-ordinary}{y}))]]}$، يكون إما~$\OLId{latin-ordinary}{x}^* \prec \OLId{latin-ordinary}{y}$ وإما~$\OLId{latin-ordinary}{x}^* = \OLId{latin-ordinary}{y}$. والحالة
  الأخيرة مستحيلة لأنها تستلزم~$\OLId{latin-ordinary}{x} \approx \OLId{latin-ordinary}{y}$، ومن ثم~$\OLId{latin-ordinary}{x}^* \prec \OLId{latin-ordinary}{y}$.
  وبالمثل، إذا كان~$\OLId{latin-ordinary}{x} \prec \OLId{latin-ordinary}{y}$ و$\OLId{latin-ordinary}{x} \not\approx \OLId{latin-ordinary}{y}$، كان
  $\OLId{latin-ordinary}{x} \prec {^*\OLId{latin-ordinary}{y}}$. لذلك، إذا كان~$\OLId{latin-ordinary}{x} \prec \OLId{latin-ordinary}{y}$ و$\OLId{latin-ordinary}{x} \not\approx \OLId{latin-ordinary}{y}$،
  كان كل~$\OLId{latin-ordinary}{w} \approx \OLId{latin-ordinary}{x}$ أصغر بحسب~$\prec$ من كل~$\OLId{latin-ordinary}{v} \approx \OLId{latin-ordinary}{y}$.
  وكل كتلة غير قياسية~$[\OLId{latin-ordinary}{x}]$ تكوّن سلسلة غير منتهية في الاتجاهين
  \[
  \cdots \prec {^{**}\OLId{latin-ordinary}{x}} \prec {^*}\OLId{latin-ordinary}{x} \prec \OLId{latin-ordinary}{x} \prec \OLId{latin-ordinary}{x}^* \prec \OLId{latin-ordinary}{x}^{**}
  \prec \cdots
  \]
  تسمى سلسلة~$\OLId{latin-looped}{Z}$ لأن نوع ترتيبها هو نوع ترتيب الأعداد الصحيحة.
\item يمكن رفع ترتيب~$\prec$ إلى الكتل: إذا كان~$\OLId{latin-ordinary}{x} \prec \OLId{latin-ordinary}{y}$ و
  $\OLId{latin-ordinary}{x} \not\approx \OLId{latin-ordinary}{y}$، كانت كتلة~$\OLId{latin-ordinary}{x}$ أصغر من كتلة~$\OLId{latin-ordinary}{y}$. وتكون الكتلة
  \emph{غير قياسية} إذا احتوت عنصرًا غير قياسي. أما الكتلة القياسية
  فهي أصغر الكتل.
\item لا توجد أصغر كتلة غير قياسية: فإذا كان~$\OLId{latin-ordinary}{y}$ غير قياسي، وُجد
  $\OLId{latin-ordinary}{x} \prec \OLId{latin-ordinary}{y}$ يكون~$\OLId{latin-ordinary}{x}$ فيه غير قياسي كذلك و$\OLId{latin-ordinary}{x} \not\approx \OLId{latin-ordinary}{y}$.
  البرهان: في النموذج القياسي~$\Struct{N}$، يقبل كل عدد القسمة على
  اثنين، وربما كان الباقي واحدًا؛ أي
  $\Sat{\OLId{latin-looped}{N}}{\lforall[\OLId{latin-ordinary}{y}][\lexists[\OLId{latin-ordinary}{x}][(\OLId{latin-ordinary}{y} = \OLId{latin-ordinary}{x} + \OLId{latin-ordinary}{x} \lor \OLId{latin-ordinary}{y} = \OLId{latin-ordinary}{x} + \OLId{latin-ordinary}{x} +
        \Obj{0}')]]}$. وبالتكافؤ الأولي، لكل~$\OLId{latin-ordinary}{y} \in \Domain{\OLId{latin-looped}{M}}$ يوجد
  $\OLId{latin-ordinary}{x} \in \Domain{\OLId{latin-looped}{M}}$ بحيث إما~$\OLId{latin-ordinary}{x} \oplus \OLId{latin-ordinary}{x} = \OLId{latin-ordinary}{y}$ وإما
  $\OLId{latin-ordinary}{x} \oplus \OLId{latin-ordinary}{x} \oplus \mathbf{z}^*= \OLId{latin-ordinary}{y}$. ولو كان~$\OLId{latin-ordinary}{x}$ قياسيًا، لكان~$\OLId{latin-ordinary}{y}$
  قياسيًا أيضًا؛ ولذلك~$\OLId{latin-ordinary}{x}$ غير قياسي، كما أن~$\OLId{latin-ordinary}{x} \prec \OLId{latin-ordinary}{y}$. وفوق ذلك،
  ينتمي~$\OLId{latin-ordinary}{x}$ و$\OLId{latin-ordinary}{y}$ إلى كتلتين مختلفتين، أي~$\OLId{latin-ordinary}{x} \not\approx \OLId{latin-ordinary}{y}$. ولرؤية
  ذلك، افترض أنهما ينتميان إلى الكتلة نفسها، أي
  $\OLId{latin-ordinary}{x} \oplus \OLId{latin-ordinary}{n} = \OLId{latin-ordinary}{y}$ لبعض~$\OLId{latin-ordinary}{n}$ قياسي. فإذا كان~$\OLId{latin-ordinary}{y} = \OLId{latin-ordinary}{x} \oplus \OLId{latin-ordinary}{x}$، كان
  $\OLId{latin-ordinary}{x} \oplus \OLId{latin-ordinary}{n} = \OLId{latin-ordinary}{x} \oplus \OLId{latin-ordinary}{x}$، ومنه~$\OLId{latin-ordinary}{x} = \OLId{latin-ordinary}{n}$ بقانون الاختزال للجمع،
  وهو قانون يصدق في~$\Struct{N}$ ومن ثم في~$\Struct{M}$ أيضًا، فيكون
  $\OLId{latin-ordinary}{x}$ قياسيًا في نهاية المطاف. وتُعالج حالة
  $\OLId{latin-ordinary}{y} = \OLId{latin-ordinary}{x} \oplus \OLId{latin-ordinary}{x} \oplus \mathbf{z}^*$ على هذا القياس.
\item وبحجة مماثلة، لا توجد أكبر كتلة.
\item ترتيب الكتل كثيف: إذا كانت~$[\OLId{latin-ordinary}{x}]$ أصغر من~$[\OLId{latin-ordinary}{y}]$، حيث
  $\OLId{latin-ordinary}{x} \not\approx \OLId{latin-ordinary}{y}$، وُجدت كتلة~$[\OLId{latin-ordinary}{z}]$ متميزة منهما وتقع بينهما. افترض
  $\OLId{latin-ordinary}{x} \prec \OLId{latin-ordinary}{y}$. وكما سبق، يقبل~$\OLId{latin-ordinary}{x} \oplus \OLId{latin-ordinary}{y}$ القسمة على اثنين، وربما
  كان ثمة باق، ولذلك يوجد~$\OLId{latin-ordinary}{u} \in \Domain{\OLId{latin-looped}{M}}$ بحيث إما
  $\OLId{latin-ordinary}{x} \oplus \OLId{latin-ordinary}{y} = \OLId{latin-ordinary}{u} \oplus \OLId{latin-ordinary}{u}$ وإما
  $\OLId{latin-ordinary}{x} \oplus \OLId{latin-ordinary}{y} = \OLId{latin-ordinary}{u} \oplus \OLId{latin-ordinary}{u} \oplus \mathbf{z}^*$. والعنصر~$\OLId{latin-ordinary}{u}$ هو
  متوسط~$\OLId{latin-ordinary}{x}$ و$\OLId{latin-ordinary}{y}$، ولذلك يقع بينهما. افترض
  $\OLId{latin-ordinary}{x} \oplus \OLId{latin-ordinary}{y} = \OLId{latin-ordinary}{u} \oplus \OLId{latin-ordinary}{u}$، والحالة الأخرى مماثلة: إذا كان
  $\OLId{latin-ordinary}{u} \approx \OLId{latin-ordinary}{x}$، فلبعض~$\OLId{latin-ordinary}{n}$ قياسي
  \[
  \OLId{latin-ordinary}{x} \oplus \OLId{latin-ordinary}{y} = \OLId{latin-ordinary}{x} \oplus \OLId{latin-ordinary}{n} \oplus \OLId{latin-ordinary}{x} \oplus \OLId{latin-ordinary}{n},
  \]
  ومن ثم~$\OLId{latin-ordinary}{y} = \OLId{latin-ordinary}{x} \oplus \OLId{latin-ordinary}{n} \oplus \OLId{latin-ordinary}{n}$، فيكون~$\OLId{latin-ordinary}{x} \approx \OLId{latin-ordinary}{y}$ خلافًا
  للفرض. فنستنتج~$\OLId{latin-ordinary}{u} \not\approx \OLId{latin-ordinary}{x}$. وتعطي حجة مماثلة
  $\OLId{latin-ordinary}{u} \not\approx \OLId{latin-ordinary}{y}$.
\end{enumerate}
وعلى هذا تُرتب الكتل غير القياسية مثل الأعداد النسبية: فهي تكوّن ترتيبًا
خطيًا~!!a{enumerable} بلا نقطتي نهاية. ويلزم من ذلك أنه لأي نموذجين
غير قياسيين~!!{enumerable}~$\Struct{M}_\OLMathNumeral{1}$ و$\Struct{M}_\OLMathNumeral{2}$ للحساب
الصادق، يكون اختزالهما إلى اللغة التي لا تحتوي إلا~$<$ و$=$ متماثلين.
فالجزآن القياسيان من~$\Struct{M_1}$ و$\Struct{M_2}$ متماثلان مع
النموذج القياسي~$\Struct{N}$، ومن ثم أحدهما مع الآخر. والكتل المكوّنة
للجزء غير القياسي مرتبة هي نفسها مثل الأعداد النسبية، ولذلك تكون
متماثلة بموجب~\olref[dlo]{thm:cantorQ}؛ ويمكن توسيع تماثل الكتل إلى
تماثل~\emph{داخل} الكتل بمقابلة عناصر اعتباطية في كل منها، ثم جعل
صورة خلف~$\OLId{latin-ordinary}{x}$ في~$\Struct{M_1}$ خلف صورة~$\OLId{latin-ordinary}{x}$ في~$\Struct{M_2}$. لاحظ
أنه لا يلزم من ذلك~\emph{أن}~$\mathfrak{M}_\OLMathNumeral{1}$ و$\mathfrak{M}_\OLMathNumeral{2}$
متماثلان في لغة الحساب الكاملة، فالتماثل دائمًا نسبي إلى توقيع؛ إذ
توجد طرائق غير متماثلة لتعريف الجمع والضرب على~$\Domain{\OLId{latin-looped}{M}_\OLMathNumeral{1}}$ و
$\Domain{\OLId{latin-looped}{M}_\OLMathNumeral{2}}$. (ويلزم ذلك كذلك من مبرهنة شهيرة لفوت، تقول إن عدد
النماذج المعدودة لنظرية تامة لا يجوز أن يكون~\OLClassicalDigits{2}.)

\begin{prob}
بيّن أنه لا يجوز أن توجد أكبر كتلة في نموذج حسابي غير قياسي.
\end{prob}

\begin{prob}
لتكن~$\Lang{L}$ !!{language} الرتبة الأولى التي لا تحتوي، عدا~$\eq$،
إلا~$<$ بوصفه~!!{predicate} وحيدًا، ولتكن
$\Struct{N} = (\Nat,<)$. جميع المجموعات الجزئية المنتهية أو المتناهية
المتممة من~$\Struct{N}$ قابلة للتعريف بغير معاملات. بيّن أنها
المجموعات الجزئية الوحيدة القابلة للتعريف بغير معاملات في
$\Struct{N}$.

(تلميح: أولًا، لتكن~$\Obj{prc}(\OLId{latin-ordinary}{x},\OLId{latin-ordinary}{y})$ صيغة~$\Lang{L}$ المختصرة لعبارة
«$\OLId{latin-ordinary}{x}$ هو السلف المباشر لـ$\OLId{latin-ordinary}{y}$»:
\[
\OLId{latin-ordinary}{x}<\OLId{latin-ordinary}{y} \land \lnot \lexists[\OLId{latin-ordinary}{z}][(\OLId{latin-ordinary}{x}<\OLId{latin-ordinary}{z} \land \OLId{latin-ordinary}{z} < \OLId{latin-ordinary}{y})].
\]
تقابل كل مجموعة جزئية من~$\Struct{N}$ قابلة للتعريف بغير معاملات
صيغة~$!A(\OLId{latin-ordinary}{x})$ في~$\Lang{L}$. ولكل صيغة كهذه~$!A$، تأمل الجملة~$!D$:
\[
\lexists[\OLId{latin-ordinary}{x}][\lforall[\OLId{latin-ordinary}{y}][\lforall[\OLId{latin-ordinary}{z}][((\OLId{latin-ordinary}{x}<\OLId{latin-ordinary}{y} \land \OLId{latin-ordinary}{x}<\OLId{latin-ordinary}{z} \land \Obj{prc}(\OLId{latin-ordinary}{y},\OLId{latin-ordinary}{z})
  \land !A(\OLId{latin-ordinary}{y})) \lif !A(\OLId{latin-ordinary}{z}))]]].
\]
بيّن أن~$\Sat{\OLId{latin-looped}{N}}{!D}$ إذا وفقط إذا كانت المجموعة الجزئية من
$\Struct{N}$ التي تعرّفها~$!A$ إما منتهية وإما متناهية المتممة.

والآن ليكن~$\Struct{M}$ نموذجًا غير قياسي مكافئًا أوليًا لـ
$\Struct{N}$. وإذا كان~$\OLId{latin-ordinary}{a} \in \Domain{\OLId{latin-looped}{M}}$ غير قياسي، فليكن
$\OLId{latin-ordinary}{b}, \OLId{latin-ordinary}{c} \in \Domain{\OLId{latin-looped}{M}}$ أكبر من~$\OLId{latin-ordinary}{a}$، وليكن~$\OLId{latin-ordinary}{b}$ السلف المباشر لـ$\OLId{latin-ordinary}{c}$.
عندئذ يوجد تشاكل ذاتي~$\OLId{latin-ordinary}{h}$ لاختزال~$\Struct{M}$ إلى لغة الترتيب بحيث
$\OLId{latin-ordinary}{h}(\OLId{latin-ordinary}{b})=\OLId{latin-ordinary}{c}$، فلماذا؟ ولذلك إذا أشبع~$\OLId{latin-ordinary}{b}$ الصيغة~$!A$، أشبعها~$\OLId{latin-ordinary}{c}$ أيضًا،
فلماذا؟ ويلزم أن~$!D$ صادقة في~$\Struct{M}$، ومن ثم في~$\Struct{N}$
كذلك. غير أن ذلك يستلزم أن المجموعة الجزئية من~$\Struct{N}$ التي
تعرّفها~$!A$ إما منتهية وإما متناهية المتممة.)
\end{prob}

\end{document}
